\chapter{Projection-Free Nonsmooth Optimization over Convex Domains}
\label{Section:projfree-domain}

\noindent\textbf{Source and author contributions.}
This chapter is based on joint work with Michael J. Neely, reported in
\cite{asgari2026projectionfree,asgari2024nonsmooth}.
I formulated the research question concerning projection-free optimization
for nonsmooth convex objectives. Neely proposed the core algorithm.
I introduced a correction to the proposed algorithm and proved its
convergence guarantees. Neely reviewed the resulting work.
I used ChatGPT, primarily GPT-6 with some use of GPT-5.6, to
refine proofs and improve the organization and editorial presentation of this
chapter.

\section{Introduction}

We study the minimization of a nonsmooth convex function over a compact
convex domain:
\begin{equation}\label{eq:pfdom:problem}
    \min_{x\in\X}f(x),\qquad
    f^*\coloneqq\min_{x\in\X}f(x).
\end{equation}
On some domains, minimizing a linear function is substantially cheaper
than computing a Euclidean projection~\cite{COMBETTES2021565,pmlr-v28-jaggi13}.
For smooth objectives, Frank--Wolfe methods avoid domain projection by
using linear optimization and a smoothness bound on the linear-model error
\cite{https://doi.org/10.1002/nav.3800030109,pmlr-v28-jaggi13}.

\subsection{Failure of direct subgradient substitution}
\label{sec:Nesterov's Counterexample}

For nonsmooth objectives, replacing the gradient in a Frank--Wolfe step
with a subgradient need not give convergence. Nesterov's counterexample
\cite[Example~1]{Nesterov2018} gives valid subgradient directions for which
the objective gap remains bounded away from zero.

\subsection{Method and scope}
\label{sec:pfdom:contributions}

The method of Asgari and Neely~\cite{asgari2026projectionfree} makes one
ordinary linear-oracle call over $\X$, one subgradient query on an auxiliary
set $\Y\supseteq\X$, and one projection onto $\Y$ per update. Here,
projection-free refers to avoiding projection onto $\X$. The method uses a
single loop, without Moreau smoothing or nested subproblems.

Under the assumptions below, let $D>0$ bound the diameter of $\X$ and
let $L>0$ bound the subgradients returned throughout $\Y$. With parameters
chosen for a query budget $T$, the returned point satisfies
$f(\bar x_T)-f^*\leq2LD/\sqrt T$
(Corollary~\ref{cor:ch2-specialization}). Both oracle counts are
$O(\epsilon^{-2})$ for objective accuracy $\epsilon$, with $L,D$ fixed.

The results below use deterministic subgradients and an exact linear
oracle. The published analysis also permits stochastic subgradients and
inexact linear optimization~\cite{asgari2024nonsmooth}. Those variants use
the same main proof ideas, with expectations and additional error terms.
We omit them so that the extra notation does not obscure the method and
its analysis. The numerical experiments are also outside this chapter's scope.

\section{Related work}
\label{sec:pfdom:context}
\label{sec:pfdom:related}

Chapter~\ref{Section:projfree} reviews methods that also handle
functional inequalities.

\subsection{From smooth conditional gradients to nonsmooth optimization}

To handle nonsmooth objectives, existing methods use richer local models,
smoothing, or dual formulations. Other methods separate feasible points
from subgradient query points.

White gives an early subgradient extension of Frank--Wolfe to
concave nondifferentiable maximization~\cite{White1993}.
Ravi et al. use approximate subdifferential sets to choose search
directions~\cite{ravi2017deterministic}. Cheung and Li construct
uniform affine approximations for separable
objectives~\cite{cheung2018solving}. The latter two methods require additional
structure or information beyond a single subgradient.

Argyriou et al. and Yurtsever et al. combine smoothing of a nonsmooth
objective term with conditional gradient steps over the
domain~\cite{argyriou2014hybrid,yurtsever2018conditional}.
Their methods require an efficiently computable proximity operator
for that term.

\subsection{Nonsmooth methods with linear and Fenchel oracles}

Lan combines randomized smoothing with the conditional gradient method
for nonsmooth Lipschitz convex objectives~\cite{lan2014complexity}.
For fixed dimension, the method attains expected objective accuracy
$\epsilon$ using $O(\epsilon^{-2})$ linear oracle calls and
$O(\epsilon^{-4})$ subgradient calls. The subgradient queries may lie
outside the original domain.

Hazan and Kale combine randomized smoothing with online Frank--Wolfe
updates for nonsmooth Lipschitz convex
losses~\cite[Theorem~4.3]{10.5555/3042573.3042808}. For losses sampled
independently from a fixed distribution, their average iterate has
a gap in expected loss of $O(T^{-1/3})$ with high probability, for fixed
dimension and confidence level. Each round uses one linear oracle
call and the gradient of the average of the smoothed losses observed
so far.

The MOLES method of Thekumparampil et al. combines Moreau smoothing
with approximate projections computed by Frank--Wolfe
steps~\cite{thekumparampil2020projection}. For nonsmooth Lipschitz
convex objectives, it attains expected objective accuracy $\epsilon$
using $O(\epsilon^{-2})$ subgradient calls and $O(\epsilon^{-2})$
linear oracle calls.
The method requires subgradient
access on a larger set containing the original domain.

Cox, Juditsky, and Nemirovski apply subgradient methods to a dual
problem obtained from a Fenchel-type representation of the
objective~\cite{cox2014dual}. Each dual subgradient is computed using
a linear oracle call over the original domain. Weights supplied by
the dual method combine the returned points into an approximate
primal solution. The approach requires the representation and
efficient proximal updates on the dual domain.

Yurtsever, Tran-Dinh, and Cevher use a Fenchel oracle that minimizes
the objective plus a linear term over the
domain~\cite{yurtsever2015universal}. This requires more than an
ordinary linear oracle. Their primal--dual framework handles
additional linear inclusion constraints through proximal updates
on the dual variables.

The domain-separation method of Asgari and Neely
\cite{asgari2026projectionfree} has the same oracle-complexity order as
MOLES, using a single loop. Section~\ref{sec:pfdom:convex} gives its updates.
Chapter~\ref{Section:projfree} extends this method to functional
inequality constraints~\cite{asgari2024nonsmooth}.

\begin{table}[htbp]
\centering
\caption{Selected approaches for general nonsmooth convex objectives.
The accuracy parameter is $\epsilon$, with dimension and other problem
parameters fixed. The counts refer to subgradient and linear oracle calls.}
\label{tab:pfdom:general-related}
\small
\renewcommand{\arraystretch}{1.25}
\setlength{\tabcolsep}{4pt}
\begin{tabular}{@{}>{\raggedright\arraybackslash}p{0.26\linewidth}>{\raggedright\arraybackslash}p{0.30\linewidth}>{\raggedright\arraybackslash}p{0.38\linewidth}@{}}
\toprule
Method & Domain access & Oracle complexity\\
\midrule
Lan's randomized smoothing~\cite{lan2014complexity}
 & Ordinary linear oracle.
 & $O(\epsilon^{-2})$ linear calls and $O(\epsilon^{-4})$ subgradient calls.\\[3pt]
MOLES~\cite{thekumparampil2020projection}
 & Ordinary linear oracle and easy auxiliary projection.
 & $O(\epsilon^{-2})$ for each oracle.\\[3pt]
Domain separation~\cite{asgari2026projectionfree}
 & Ordinary linear oracle and easy auxiliary projection.
 & $O(\epsilon^{-2})$ for each oracle.\\
\bottomrule
\end{tabular}
\end{table}

For general Lipschitz convex objectives, the classical worst-case
lower bound is $\Omega(\epsilon^{-2})$ first-order oracle calls
to reach objective accuracy $\epsilon$, provided the dimension is
sufficiently large~\cite{nemirovskij1983problem,bubeck2015convex}.
MOLES and the method of Asgari and Neely match this order in their
subgradient counts under their respective assumptions.

Lan proves that $\Omega(\epsilon^{-2})$ linear oracle calls are
necessary in the worst case for general Lipschitz convex objectives
in sufficiently high dimension~\cite[Section~2.3]{lan2014complexity}.
His model requires the output to be a convex combination of the
initial point and the points returned by the linear oracle.
His randomized smoothing method attains the optimal order in
$\epsilon$ for linear oracle calls, with an additional dimension
factor, and uses $O(\epsilon^{-4})$ subgradient calls.
MOLES and our general convex method have the required output form
and match the linear oracle lower bound without the additional dimension
factor, using $O(\epsilon^{-2})$ calls to each oracle.


For the domain-separation method, the bound
$f(\bar x_T)-f^*\leq2LD/\sqrt T$ gives at most
$\lceil4L^2D^2/\epsilon^2\rceil$ linear-oracle calls, including initialization,
to attain objective accuracy $\epsilon$, for $L,D,\epsilon>0$.
Its leading constant is smaller than the
corresponding worst-case bound in the MOLES analysis
\cite[Appendix~C.2.2]{thekumparampil2020projection}.

\section{Problem and oracle assumptions}
\label{sec:pfdom:setup}

We follow the notation of Section~\ref{sec:intro:notation}.
We work in $\Real^n$ with the Euclidean inner product and norm $\|\cdot\|$.
The function $f:\Real^n\to\Real$ is finite-valued and convex, and
$\X\subseteq\Y\subseteq\Real^n$.
The vectors $x_t,y_t,q_t\in\Real^n$ denote the domain point, auxiliary query point,
and equality multiplier. Throughout this chapter, the integer $T\geq1$ is a query
budget for each oracle: the algorithm uses at most $T$ calls to each oracle,
including any linear-oracle call used for initialization.

\Needspace{7\baselineskip}
\begin{Thm:Assumption}[Domain and ordinary oracles]\label{ass:pfdom:oracles}
The following conditions hold.
\begin{enumerate}[label=(\roman*)]
\item The set $\X$ is nonempty, compact, and convex, with Euclidean diameter
\[
    \operatorname{diam}(\X)\coloneqq\sup_{x,z\in\X}\|x-z\|\leq D,
\]
where $0\leq D<\infty$ is a known upper bound, not necessarily the exact diameter.
The closed convex set $\Y$ contains $\X$ and has an
available Euclidean projection $\Pi_\Y$.
\item A linear optimization oracle is available. For each direction
$v\in\Real^n$, it returns
\[
    \operatorname{LOO}_\X(v)\in\argmin_{x\in\X}\inner{v}{x}.
\]
\item At every query $y\in\Y$, the subgradient oracle returns $s$ with
\[
    s\in\partial f(y),\qquad \|s\|\leq L,
\]
where $0\leq L<\infty$ is known.
\end{enumerate}
\end{Thm:Assumption}

An initial point can be obtained by calling $\operatorname{LOO}_\X(0)$,
or any available point in $\X$ can be used. The numerical values of $L$
and $D$ are used to choose the algorithm parameters and evaluate the
convergence guarantee.

Lemma~\ref{lem:intro:lipschitz}, with the subgradient norm bound
in Assumption~\ref{ass:pfdom:oracles}, gives
\begin{equation}\label{eq:pfdom:lipschitz}
    |f(y)-f(w)|\leq L\|y-w\|,\qquad y,w\in\Y.
\end{equation}
In particular, $f$ is continuous on $\X$. Since $\X$ is nonempty and
compact, the Weierstrass theorem guarantees a minimizer $x^*\in\X$
\cite{bertsekas1997nonlinear}.

\subsection{Choosing the auxiliary set}
\label{sec:pfdom:auxiliary-choice}

The set $\Y$ must contain $\X$, admit an available Euclidean projection,
and support the required bounded subgradient queries. It need not be bounded:
if the subgradient bound holds throughout $\Real^n$, we can take $\Y\coloneqq\Real^n$ and the
projection is the identity.
Given $x_1\in\X$ and the diameter bound $D$, two bounded choices are
\begin{equation}\label{eq:pfdom:enclosing-sets}
 \Y\coloneqq\{y:\|y-x_1\|\leq D\}
 \quad\text{or}\quad
 \Y\coloneqq\{y:|y_i-x_{1,i}|\leq D\ \text{for all }i\}.
\end{equation}
Both contain $\X$: for $x\in\X$, Assumption~\ref{ass:pfdom:oracles}(i)
gives $|x_i-x_{1,i}|\leq\|x-x_1\|\leq D$ for every coordinate $i$.
Projection onto the closed ball rescales a displacement
that exceeds its radius, while projection onto the box clips each coordinate
\cite[Section~3.2]{COMBETTES2021565}.

A subgradient bound on $\X$ alone is not enough; it must hold on the
chosen $\Y$. In the guarantee $f(\bar x_T)-f^*\leq2LD/\sqrt T$ of
Corollary~\ref{cor:ch2-specialization}, $D$ bounds the diameter of $\X$.
For fixed $\X$ and $T$, enlarging $\Y$ therefore weakens this guarantee
only if a larger uniform bound $L$ is needed.
Section~\ref{sec:pfdom:auxiliary-example} gives an example and a
Lipschitz extension. An extension's existence does not supply an inexpensive
subgradient oracle: both the bound and oracle access on $\Y$ are required.

\section{Derivation and algorithm}
\label{sec:pfdom:convex}
\label{sec:pfdom:separation}

The failure in Section~\ref{sec:Nesterov's Counterexample} motivates a
different linear-oracle direction. Here the oracle uses a multiplier
that accumulates disagreement between feasible points and auxiliary
query points.

Introducing $y=x$ gives the equivalent problem
\[
    \min_{x\in\X,\,y\in\Y}\{f(y):y=x\}.
\]
The Lagrangian is
\[
    \mathcal L(x,y,q)\coloneqq f(y)+\inner{q}{x-y}.
\]
For a fixed multiplier $q$, minimizing the Lagrangian separates into
minimizing $\inner{q}{x}$ over $\X$ and $f(y)-\inner{q}{y}$ over $\Y$.
Even at an optimal multiplier, these subproblems can have multiple
minimizers, so independently selected minimizers need not satisfy
$x=y$~\cite[Section~5.5.5]{boyd_vandenberghe_2004}.
The multiplier update below gives the consensus identity
\eqref{eq:pfdom:consensus}.
The linear oracle chooses $x_{t+1}$ by minimizing its
$x$-dependent term at $q_t$. At iteration $t$, let $s_t$ denote the
subgradient returned by the oracle at $y_t$, satisfying
Assumption~\ref{ass:pfdom:oracles}(iii). The relevant subgradient and derivative are
\[
    \partial_y\mathcal L(x,y,q)=\partial f(y)-q,
    \qquad
    \nabla_q\mathcal L(x,y,q)=x-y.
\]
These expressions suggest descent in $y$ and ascent in $q$.
For $\alpha,\eta>0$, take a projected subgradient step in $y$ of size
$1/\alpha$, using the updated multiplier $q_{t+1}$, and a gradient ascent
step in $q$ of size $\eta$, using the new points $x_{t+1}$ and $y_{t+1}$.
This gives the coupled updates
\begin{align}\label{eq:ch2-implicit-pair}
    y_{t+1}&=\Pi_\Y\left(y_t+\frac{q_{t+1}-s_t}{\alpha}\right),&
    q_{t+1}&=q_t-\eta(y_{t+1}-x_{t+1}).
\end{align}
Eliminating $q_{t+1}$ gives an explicit update with one projection.
For a nonempty closed convex $\Y$, first-order optimality for
$\min_{w\in\Y}\|w-v\|^2/2$ characterizes $p=\Pi_\Y(v)$ by
$\inner{p-v}{w-p}\geq0$ for every $w\in\Y$.
This is the condition in Lemma~\ref{lem:intro:projection} with $Q=\Y$.
Taking $p=y_{t+1}$ and $v=y_t+(q_{t+1}-s_t)/\alpha$ in this
condition and multiplying by $\alpha>0$ gives
\begin{equation}\label{eq:ch2-basic-subproblem}
    \inner{s_t-q_{t+1}+\alpha(y_{t+1}-y_t)}{w-y_{t+1}}
    \geq0\qquad(w\in\Y).
\end{equation}
Substituting the multiplier update~\eqref{eq:ch2-implicit-pair} into
the first argument of the inner product in~\eqref{eq:ch2-basic-subproblem} gives
\[
    (\alpha+\eta)y_{t+1}
      -(\alpha y_t+\eta x_{t+1}+q_t-s_t).
\]
Dividing the resulting condition by $\alpha+\eta>0$ and applying the
same projection characterization gives the explicit update
\begin{equation}\label{eq:ch2-basic-updates}
    y_{t+1}\coloneqq\Pi_\Y\left(
    \frac{\alpha y_t+\eta x_{t+1}+q_t-s_t}{\alpha+\eta}\right).
\end{equation}

The linear oracle in Assumption~\ref{ass:pfdom:oracles}(ii) uses direction
$q_t$. Comparing its minimizer $x_{t+1}$ with the feasible point $x^*\in\X$ gives
\begin{equation}\label{eq:pfdom:loo-optimality}
    \inner{q_t}{x_{t+1}-x^*}\leq0.
\end{equation}

\begin{algorithm}[H]
\caption{Domain separation for a nonsmooth convex objective}
\label{alg:pfdom:convex}
\begin{algorithmic}[1]
\Require Integer query budget $T\geq1$, $\alpha,\eta>0$, $x_1\in\X$
\State $y_1\gets x_1$, $q_1\gets0$
\For{$t=1,\ldots,T-1$}
    \State $x_{t+1}\gets\operatorname{LOO}_\X(q_t)$
    \State Query $s_t$ at $y_t$
    \State $\displaystyle y_{t+1}\gets\Pi_\Y\left(
      \frac{\alpha y_t+\eta x_{t+1}+q_t-s_t}{\alpha+\eta}\right)$
    \State $q_{t+1}\gets q_t-\eta(y_{t+1}-x_{t+1})$
\EndFor
\State \Return $\displaystyle\bar x_T\coloneqq\frac1T\sum_{t=1}^T x_t$
\end{algorithmic}
\end{algorithm}

The algorithm performs $T-1$ updates. Each update uses one ordinary linear-oracle
call, one subgradient call, and one projection onto $\Y$.

\section{Convergence analysis}
\label{sec:pfdom:convergence}

\subsection{Objective bound and parameter choice}

\begin{Thm:Theorem}[General convex objective bound]\label{thm:pfdom:convex}
Under Assumption~\ref{ass:pfdom:oracles}, Algorithm~\ref{alg:pfdom:convex}
with any positive $\alpha,\eta$ and integer $T\geq1$ satisfies
\begin{equation}\label{eq:ch2-basic-bound}
    0\leq f(\bar x_T)-f^*
    \leq\frac{L^2}{2\alpha}+\frac{\alpha D^2}{2T}
      +\frac{(T-1)\eta D^2}{2T}
      +\frac{L^2}{2\eta T}.
\end{equation}
\end{Thm:Theorem}

\Needspace{13\baselineskip}
\begin{Thm:Corollary}[Convex convergence rate with explicit parameter choices]\label{cor:ch2-specialization}
Under Assumption~\ref{ass:pfdom:oracles}, assume $L,D>0$, fix an integer
$T\geq1$, and take
\begin{equation}\label{eq:pfdom:tuned-parameters}
    \eta\coloneqq\frac{L}{D\sqrt T},\qquad
    \alpha\coloneqq T\eta.
\end{equation}
Then Algorithm~\ref{alg:pfdom:convex} satisfies
\begin{equation}\label{eq:ch2-basic-rate}
    f(\bar x_T)-f^*
    \leq\frac{2LD}{\sqrt T}.
\end{equation}
\end{Thm:Corollary}
\begin{proof}
The choices in~\eqref{eq:pfdom:tuned-parameters} are positive. Substituting
them into~\eqref{eq:ch2-basic-bound} gives
\[
 f(\bar x_T)-f^*
 \le\frac{LD}{\sqrt T}
       +\left(\frac{T-1}{T}+1\right)
        \frac{LD}{2\sqrt T}
 \le\frac{2LD}{\sqrt T}.
\]
\end{proof}

The first-order and ordinary linear-oracle complexity is $O(\epsilon^{-2})$.
Indeed, for $L,D,\epsilon>0$, choose
\begin{equation}\label{eq:pfdom:accuracy-budget}
 T\coloneqq\left\lceil\left(\frac{2LD}{\epsilon}\right)^2\right\rceil.
\end{equation}
The ceiling ensures $T\geq1$ and $\sqrt T\geq2LD/\epsilon$, so
\eqref{eq:ch2-basic-rate} gives $f(\bar x_T)-f^*\leq\epsilon$.
Algorithm~\ref{alg:pfdom:convex} makes $T-1$ calls to each oracle,
with at most one additional linear-oracle call to initialize $x_1$.
Thus both oracle counts respect the query budget $T$.
Since $T\leq(2LD/\epsilon)^2+1$ in~\eqref{eq:pfdom:accuracy-budget},
these counts have the claimed order for fixed $L,D$ as $\epsilon\downarrow0$.

\subsection{Proof of the objective bound}

The following lemma converts a cumulative bound at the query points
into an objective bound at $\bar x_T$. Lipschitz continuity bounds the
objective error caused by disagreement between the two averages, and
the quadratic multiplier term controls this error.

\begin{Thm:Lemma}[Objective bound at the returned point]
\label{lem:pfdom:equality-transfer}
Let $T\geq1$ be an integer, $\eta>0$, and let $x_t,y_t$ be
points in a convex set $\Y\subseteq\Real^n$.
Set $\bar x_T\coloneqq T^{-1}\sum_{t=1}^T x_t$, initialize $q_1\coloneqq0$ and $x_1=y_1$,
and update
\begin{equation}\label{eq:pfdom:transfer-multiplier-update}
    q_{t+1}\coloneqq q_t-\eta(y_{t+1}-x_{t+1}),\qquad t=1,\ldots,T-1.
\end{equation}
Then the following identities hold:
\begin{align}
    -\inner{q_t}{y_{t+1}-x_{t+1}}
      +\frac\eta2\|y_{t+1}-x_{t+1}\|^2
    &=\frac{\|q_{t+1}\|^2-\|q_t\|^2}{2\eta},
      \label{eq:lem2_q_squared}\\
    q_T&=\eta T(\bar x_T-\bar y_T),\qquad
    \bar y_T\coloneqq\frac1T\sum_{t=1}^T y_t.
      \label{eq:pfdom:consensus}
\end{align}
Suppose additionally that $f:\Y\to\Real$ is convex and $L$-Lipschitz,
with finite $L$. For any real reference value $f_c$ and finite constant $C$,
\begin{equation}\label{eq:pfdom:transfer-potential}
    \sum_{t=1}^T[f(y_t)-f_c]+\frac{\|q_T\|^2}{2\eta}\leq C
\end{equation}
implies
\begin{equation}\label{eq:pfdom:objective-transfer}
    f(\bar x_T)-f_c\leq\frac CT+\frac{L^2}{2\eta T}.
\end{equation}
\end{Thm:Lemma}
\begin{proof}
For $t=1,\ldots,T-1$, substituting the multiplier update
\eqref{eq:pfdom:transfer-multiplier-update} into $\|q_{t+1}\|^2$ and
expanding gives
\begin{align*}
 \|q_{t+1}\|^2
 &=\|q_t-\eta(y_{t+1}-x_{t+1})\|^2\\
 &=\|q_t\|^2-2\eta\inner{q_t}{y_{t+1}-x_{t+1}}
   +\eta^2\|y_{t+1}-x_{t+1}\|^2.
\end{align*}
Subtracting $\|q_t\|^2$ from both sides and dividing by $2\eta>0$
yields~\eqref{eq:lem2_q_squared}.

Summing~\eqref{eq:pfdom:transfer-multiplier-update} and using $q_1=0$
and $x_1=y_1$ gives
\begin{align*}
 q_T
 &=q_T-q_1
 =\sum_{t=1}^{T-1}(q_{t+1}-q_t)
 =-\eta\sum_{t=1}^{T-1}(y_{t+1}-x_{t+1})\\
 &=\eta\sum_{t=2}^{T}(x_t-y_t)
 =\eta\sum_{t=1}^{T}(x_t-y_t)
 =\eta T(\bar x_T-\bar y_T).
\end{align*}
The first line telescopes; the second shifts the index and inserts the
zero term $x_1-y_1$, proving~\eqref{eq:pfdom:consensus}.

Since $\Y$ is convex, both averages belong to $\Y$, and
\begin{align*}
 T\{f(\bar x_T)-f_c\}
 &\overset{\text{(A)}}{\le}T\{f(\bar y_T)-f_c\}+LT\|\bar x_T-\bar y_T\|\\
 &\overset{\text{(B)}}{\le}\sum_{t=1}^T\{f(y_t)-f_c\}+\frac L\eta\|q_T\|.
\end{align*}
Here (A) uses the lemma's Lipschitz assumption at $\bar x_T,\bar y_T$;
(B) uses Jensen's inequality~\eqref{eq:intro:jensen} with points $y_t$
and weights $1/T$, followed by~\eqref{eq:pfdom:consensus} and $\eta>0$.
Using~\eqref{eq:pfdom:transfer-potential} and completing the square gives
\begin{align*}
 T\{f(\bar x_T)-f_c\}
 &\le C+\frac L\eta\|q_T\|-
                 \frac{\|q_T\|^2}{2\eta}\\
 &=C+\frac{L^2}{2\eta}
   -\frac1{2\eta}(\|q_T\|-L)^2\\
 &\le C+\frac{L^2}{2\eta}.
\end{align*}
Dividing by $T$ proves~\eqref{eq:pfdom:objective-transfer}.
\end{proof}

\begin{proof}[Proof of Theorem~\ref{thm:pfdom:convex}]
\par\medskip\noindent\textbf{Primal comparison.}
Fix a minimizer $x^*\in\X$. For $t=1,\ldots,T-1$, define the convex
function
\[
 \psi_t(y)\coloneqq\inner{s_t-q_t}{y}+\frac\eta2\|y-x_{t+1}\|^2.
\]
The explicit projection update~\eqref{eq:ch2-basic-updates} is equivalent to
\begin{equation}\label{eq:pfdom:quadratic-minimization}
 y_{t+1}=\argmin_{y\in\Y}
 \left\{\psi_t(y)+\frac\alpha2\|y-y_t\|^2\right\}.
\end{equation}
Indeed, expanding the objective in~\eqref{eq:pfdom:quadratic-minimization}
gives
\begin{align*}
 \psi_t(y)+\frac\alpha2\|y-y_t\|^2
 &=\frac{\alpha+\eta}{2}\|y\|^2
   -\inner{\alpha y_t+\eta x_{t+1}+q_t-s_t}{y}\\
 &\quad+\frac\alpha2\|y_t\|^2+\frac\eta2\|x_{t+1}\|^2.
\end{align*}
Completing the square therefore gives the projection
in~\eqref{eq:ch2-basic-updates}.

Apply~\eqref{eq:pushback} with $f=\psi_t$, coefficient $\alpha$,
reference point $y_t$, comparison point $x^*\in\Y$, and quadratic geometry
\[
 C=\Real^n,\qquad Q=\Y,\qquad
 \Phi(y)\coloneqq\frac12\|y\|^2,\qquad
 D_\Phi(u,v)=\frac12\|u-v\|^2.
\]
Here $Q_{\mathrm{ref}}=\Y$, so $y_t,y_{t+1}$ are admissible reference points.
The three-point inequality gives (A):
\begin{align*}
 &\psi_t(y_{t+1})+\frac\alpha2\|y_{t+1}-y_t\|^2\\
 &\quad\overset{\mathrm{(A)}}{\le}
 \psi_t(x^*)+\frac\alpha2\|x^*-y_t\|^2
   -\frac\alpha2\|x^*-y_{t+1}\|^2.
\end{align*}
Expanding $\psi_t$ and moving its linear terms to the left gives
\begin{equation}\label{eq:pfdom:quadratic-comparison}
\begin{split}
 &\inner{s_t-q_t}{y_{t+1}-x^*}
   +\frac\eta2\|y_{t+1}-x_{t+1}\|^2
   +\frac\alpha2\|y_{t+1}-y_t\|^2\\
 &\quad\le\frac\eta2\|x_{t+1}-x^*\|^2
   +\frac\alpha2\bigl(\|y_t-x^*\|^2-\|y_{t+1}-x^*\|^2\bigr).
\end{split}
\end{equation}

\par\medskip\noindent\textbf{Multiplier and subgradient terms.}
Splitting $y_{t+1}-x^*=(y_{t+1}-x_{t+1})+(x_{t+1}-x^*)$
in~\eqref{eq:pfdom:quadratic-comparison} gives
\begin{equation}\label{eq:pfdom:multiplier-comparison}
\begin{split}
 &-\inner{q_t}{y_{t+1}-x^*}
   +\frac\eta2\|y_{t+1}-x_{t+1}\|^2\\
 &\quad=-\inner{q_t}{y_{t+1}-x_{t+1}}
   +\frac\eta2\|y_{t+1}-x_{t+1}\|^2
   -\inner{q_t}{x_{t+1}-x^*}\\
 &\quad\overset{\mathrm{(A)}}{=}
   \frac{\|q_{t+1}\|^2-\|q_t\|^2}{2\eta}
   -\inner{q_t}{x_{t+1}-x^*}\\
 &\quad\overset{\mathrm{(B)}}{\ge}
   \frac{\|q_{t+1}\|^2-\|q_t\|^2}{2\eta}.
\end{split}
\end{equation}
Here (A) uses the multiplier identity~\eqref{eq:lem2_q_squared}, and
(B) uses the linear-oracle inequality~\eqref{eq:pfdom:loo-optimality}.

To handle the subgradient term, write
\begin{equation}\label{eq:pfdom:subgradient-young}
\begin{aligned}
 f(y_t)-f^*
 &\overset{\mathrm{(A)}}{\le}\inner{s_t}{y_t-x^*}\\
 &=\inner{s_t}{y_{t+1}-x^*}+\inner{s_t}{y_t-y_{t+1}}\\
 &\overset{\mathrm{(B)}}{\le}\inner{s_t}{y_{t+1}-x^*}
   +\frac{\|s_t\|^2}{2\alpha}
   +\frac\alpha2\|y_{t+1}-y_t\|^2\\
 &\overset{\mathrm{(C)}}{\le}\inner{s_t}{y_{t+1}-x^*}
   +\frac{L^2}{2\alpha}
   +\frac\alpha2\|y_{t+1}-y_t\|^2.
\end{aligned}
\end{equation}
Here (A) is the subgradient inequality~\eqref{eq:intro:subgradient}
at $y_t$, with subgradient $s_t$ and comparison point $x^*$,
using $f(x^*)=f^*$. Step (B) uses Young's
inequality~\eqref{eq:intro:holder-young} with
\[
 g=s_t,\qquad v=y_t-y_{t+1},\qquad\gamma=\alpha>0.
\]
Step (C) uses $\|s_t\|\le L$ from
Assumption~\ref{ass:pfdom:oracles}(iii).

Substituting~\eqref{eq:pfdom:multiplier-comparison}
and~\eqref{eq:pfdom:subgradient-young}
into~\eqref{eq:pfdom:quadratic-comparison} now gives
\begin{equation}\label{eq:pfdom:one-step}
\begin{split}
 &f(y_t)-f^*
   +\frac{\|q_{t+1}\|^2-\|q_t\|^2}{2\eta}\\
 &\quad\overset{\mathrm{(A)}}{\le}
   \frac\alpha2\bigl(\|y_t-x^*\|^2-\|y_{t+1}-x^*\|^2\bigr)
   +\frac\eta2\|x_{t+1}-x^*\|^2+\frac{L^2}{2\alpha}\\
 &\quad\overset{\mathrm{(B)}}{\le}
   \frac\alpha2\bigl(\|y_t-x^*\|^2-\|y_{t+1}-x^*\|^2\bigr)
   +\frac{\eta D^2}{2}+\frac{L^2}{2\alpha}.
\end{split}
\end{equation}
For (A), \eqref{eq:pfdom:multiplier-comparison} bounds the multiplier
terms from below, while~\eqref{eq:pfdom:subgradient-young} bounds the
remaining subgradient and movement terms from below by
$f(y_t)-f^*-L^2/(2\alpha)$. Step (B) uses $x_{t+1},x^*\in\X$
and the diameter bound in Assumption~\ref{ass:pfdom:oracles}(i).

\par\medskip\noindent\textbf{Summation and terminal term.}
Summing~\eqref{eq:pfdom:one-step} over $t=1,\ldots,T-1$, the two
differences telescope:
\begin{align*}
 \sum_{t=1}^{T-1}\bigl(\|q_{t+1}\|^2-\|q_t\|^2\bigr)
 &=\|q_T\|^2-\|q_1\|^2,\\
 \sum_{t=1}^{T-1}\bigl(\|y_t-x^*\|^2-\|y_{t+1}-x^*\|^2\bigr)
 &=\|y_1-x^*\|^2-\|y_T-x^*\|^2.
\end{align*}
Consequently, using $q_1=0$,
\begin{equation}\label{eq:pfdom:summed-one-step}
\begin{split}
 &\sum_{t=1}^{T-1}[f(y_t)-f^*]+\frac{\|q_T\|^2}{2\eta}\\
 &\quad\le\frac\alpha2\|y_1-x^*\|^2
   -\frac\alpha2\|y_T-x^*\|^2
   +(T-1)\left(\frac{L^2}{2\alpha}+\frac{\eta D^2}{2}\right)\\
 &\quad\overset{\mathrm{(A)}}{\le}\frac{\alpha D^2}{2}
   -\frac\alpha2\|y_T-x^*\|^2
   +(T-1)\left(\frac{L^2}{2\alpha}+\frac{\eta D^2}{2}\right).
\end{split}
\end{equation}
Here (A) uses $y_1=x_1\in\X$, $x^*\in\X$, and
Assumption~\ref{ass:pfdom:oracles}(i).

The sum in~\eqref{eq:pfdom:summed-one-step} does not yet include
$f(y_T)-f^*$. For this terminal term,
\begin{equation}\label{eq:pfdom:terminal}
\begin{aligned}
 f(y_T)-f^*&\overset{\mathrm{(A)}}{\le}L\|y_T-x^*\|\\
 &\overset{\mathrm{(B)}}{\le}\frac{L^2}{2\alpha}
   +\frac\alpha2\|y_T-x^*\|^2.
\end{aligned}
\end{equation}
Step (A) uses Lipschitz continuity~\eqref{eq:pfdom:lipschitz}
at $y_T,x^*\in\Y$. Step (B) uses the scalar form of Young's
inequality~\eqref{eq:intro:holder-young} with $g=L$,
$v=\|y_T-x^*\|$, and $\gamma=\alpha>0$.

Adding~\eqref{eq:pfdom:terminal} to~\eqref{eq:pfdom:summed-one-step}
cancels the terminal squared-distance terms. Thus
\begin{equation}\label{eq:pfdom:potential-bound}
\begin{split}
 &\sum_{t=1}^T[f(y_t)-f^*]+\frac{\|q_T\|^2}{2\eta}\\
 &\quad\le\frac{\alpha D^2}{2}
   +(T-1)\left(\frac{L^2}{2\alpha}+\frac{\eta D^2}{2}\right)
   +\frac{L^2}{2\alpha}\\
 &\quad=\frac{TL^2}{2\alpha}+\frac{\alpha D^2}{2}
   +\frac{(T-1)\eta D^2}{2}\eqqcolon C_T^{(0)}.
\end{split}
\end{equation}
For $T=1$, the sums over $t=1,\ldots,T-1$ are empty, and the same
calculation remains valid.

\par\medskip\noindent\textbf{Output bound.}
Apply Lemma~\ref{lem:pfdom:equality-transfer} with $f_c=f^*$
and $C=C_T^{(0)}$. Its points lie in $\Y$ because
$x_t\in\X\subseteq\Y$ and $y_t\in\Y$. Its initialization and
multiplier-update hypotheses hold by Algorithm~\ref{alg:pfdom:convex},
and its convexity and Lipschitz hypotheses follow from the convexity of $f$
and~\eqref{eq:pfdom:lipschitz}. Therefore,
\begin{equation}\label{eq:pfdom:convex-transfer}
\begin{aligned}
 f(\bar x_T)-f^*
 &\overset{\mathrm{(A)}}{\le}\frac{C_T^{(0)}}{T}
   +\frac{L^2}{2\eta T}\\
 &=\frac1T\left(\frac{TL^2}{2\alpha}+\frac{\alpha D^2}{2}
   +\frac{(T-1)\eta D^2}{2}\right)+\frac{L^2}{2\eta T}\\
 &=\frac{L^2}{2\alpha}+\frac{\alpha D^2}{2T}
   +\frac{(T-1)\eta D^2}{2T}+\frac{L^2}{2\eta T}.
\end{aligned}
\end{equation}
Here (A) is precisely the objective-transfer
inequality~\eqref{eq:pfdom:objective-transfer},
using~\eqref{eq:pfdom:potential-bound}.
Since every $x_t\in\X$ and $\X$ is convex, $\bar x_T\in\X$.
Hence $f(\bar x_T)\ge f^*$, proving the lower bound
in~\eqref{eq:ch2-basic-bound} as well.
\end{proof}

\subsection{Subgradient bounds on auxiliary sets}
\label{sec:pfdom:auxiliary-example}

Consider
\[
 f(x)\coloneqq\max\{e^{-x},e^x\},\qquad \X\coloneqq[-1,1].
\]
This function has subgradients bounded by $e$ on $\X$, but has no finite
global subgradient bound. If $x_1\coloneqq1/2$ and $D\coloneqq2$, the enclosing ball in
\eqref{eq:pfdom:enclosing-sets} is $\Y=[-3/2,5/2]$, where the required
bound increases to $L=e^{5/2}$. Indeed, $f(x)=e^{|x|}$ has derivative
magnitude $e^{|x|}$ for $x\ne0$ and subdifferential $[-1,1]$ at zero;
the largest $|x|$ in this auxiliary interval is $5/2$.

For a convex $L$-Lipschitz function on the convex set $\X$, the
construction
\begin{equation}\label{eq:pfdom:lipschitz-extension}
 \widetilde f(y)\coloneqq\inf_{z\in\X}\{f(z)+L\|y-z\|\}
\end{equation}
is convex and $L$-Lipschitz on $\Real^n$ and agrees with $f$ on $\X$
by the McShane--Whitney extension theorem
\cite[Theorem~2.3 and Proposition~3.2]{lmcs:6105}; see also
\cite[Appendix~D]{asgari2024nonsmooth} for its use in projection-free
optimization. In the preceding example it gives
$\widetilde f(x)=f(x)$ for $|x|\leq1$ and $\widetilde f(x)=e|x|$ otherwise.
For the latter formula, take $L\coloneqq e$ in~\eqref{eq:pfdom:lipschitz-extension}.
If $x>1$ and $z\in[-1,1]$, the subgradient inequality
\eqref{eq:intro:subgradient} at $1$ with subgradient $e$ gives
$f(z)\geq e+e(z-1)=ez$. Hence
$f(z)+e|x-z|\geq ex$, with equality at $z=1$; symmetry gives the case $x<-1$.
The extension preserves the values on $\X$, but evaluating it involves
optimization over $\X$; its existence does not provide an inexpensive
subgradient oracle.

\section{Discussion and conclusion}
\label{sec:pfdom:discussion}

The method returns a feasible average while evaluating subgradients
throughout the auxiliary set. Its implementation is useful when projection
onto this set and linear optimization over $\X$ are inexpensive.
The bound depends on a uniform subgradient norm bound on the auxiliary
set, and the stated parameter choice uses the prescribed query budget.

Chapter~\ref{Section:projfree} extends this construction to functional
inequalities. It keeps the ordinary linear oracle for $\X$ and uses
additional multipliers for those constraints.
